% Continuation of the proof of Theorem 2.1.
One easily concludes that $(\mathfrak g/\mathfrak h)^T=
(\mathfrak g/\mathfrak h)^{\ad(a)}$, and since by hypothesis the left-hand
side is zero, one concludes (vii).

Let us prove (vii) $\Rightarrow$ (vi). This implication is contained in
the following result, which makes 2.1 more precise:
\begin{corollary}{2.2}\label{XIII.2.2}
Let $G$ be a smooth algebraic group over a field $k$, $H$ a smooth
algebraic subgroup, $N$ its normalizer in $G$, $\varphi:G\times H\to G$
the morphism defined by $\varphi(g,h)=\ad(g)h=ghg^{-1}$,
$\psi:X=G\times^N H\to G$ the morphism deduced from $\varphi$ by passage
to the quotient (cf. \No 1), and $a\in H(k)$. The following conditions
are equivalent:
\begin{enumeratei}
\item $\varphi$ is smooth at $(e,a)$.
\item $\psi$ is \'etale at $(\bar e,a)$, and $N$ is smooth over $k$.
\item $(\mathfrak g/\mathfrak h)^{\ad(a)}=0$ (where $\mathfrak g,\mathfrak h$
are the Lie algebras of $G,H$).
\end{enumeratei}
These conditions imply $H^0=N^0$.
\end{corollary}
One knows that the smoothness of $\varphi$ (which is a morphism of smooth
$k$-preschemes) at a point rational over $k$ is equivalent to the
surjectivity of the tangent map at that point. Now an immediate
calculation shows that this tangent map is written (by means of the
usual identifications of the tangent spaces at points of $G$ and $H$
with the Lie algebra of $G$ and $H$)
\[
 d\varphi(\xi,\eta)=(\mathrm{id}-\ad(a))\cdot\xi+\eta,
\]
considered as a map from $\mathfrak g\times\mathfrak h$ to $\mathfrak g$.
Surjectivity is therefore equivalent to the surjectivity of
$(\mathrm{id}-\ad(a))$ on $\mathfrak g/\mathfrak h$, i.e. to (iii).
Now (iii) implies a fortiori $(\mathfrak g/\mathfrak h)^H=0$, i.e.
(cf. 2.0, where the connectedness hypothesis at the beginning of the
section is unnecessary) $H^0=N^0$. One deduces that $N$ is smooth, and
$\dim H=\dim N$, whence $\dim X=\dim G$. Now, since $q:G\times H\to X$
is flat and $\psi=\varphi\circ q$ smooth at $(e,a)$, it follows that
$\psi$ is smooth at $q(e,a)=(\bar e,a)$, hence \'etale at this point
for reasons of dimension. Thus one has proved (i) $\Leftrightarrow$
(iii) $\Rightarrow$ (ii); on the other hand (ii) $\Rightarrow$ (i),
since the smoothness of $N$ implies that of $q$. \hfill Q.E.D.
% typo?: kept as printed: psi=varphi circ q in this proof.

Finally, let us prove (vi) $\Rightarrow$ (i), which, together with the
implications already established, will prove the theorem. Suppose first
that $G$ is affine. Let $U$ be a nonempty open subset of $G$ such that
$x\in U(k)$ implies that $x$ is contained in a conjugate of $H$. Let $C$
be a Cartan subgroup of $G$. Using the implication (i) $\Rightarrow$ (iv)
for $C$ in place of $H$ (this is Borel's ``density theorem''), it follows
that one can find a conjugate of $C$ which meets $U$; hence one may
suppose that $U\cap C\ne\varnothing$, i.e. that there exists a nonempty
open subset $V$ of $C$ such that for every $x\in V(k)$, $x$ is contained
in a conjugate of $H$. Write $C$ as a product
\[
 C=T\cdot C_u
\]
where $T$ is the maximal torus of $C$ (which is a maximal torus of $G$)
and $C_u$ the unipotent part of $C$, $T$ being in the center of $C$
(BIBLE 6 th. 2). One may again suppose that $k$ has infinite transcendence
degree over its prime subfield, which allows one to find an element $t$
of $T(k)$ belonging to the projection of $V$ onto $T$ (which is a nonempty
open subset of $T$), i.e. $t\cdot C_u\cap V\ne\varnothing$, and such
that $t$ ``generates'' $T$. Since every algebraic subgroup of $G$ which
contains a product $t\cdot u$ ($t\in T(k)$, $u\in C_u(k)$) contains both
factors (BIBLE 4 th. 3), it follows, with the preceding choice of $t$,
and taking $t\cdot u\in V(k)$, that there exists a conjugate of $H$ which
contains $t$, hence $T$. Thus one may already suppose that one has
\[
 T\subset H.
\]
% original p. 260
If $W$ is the open subset of $C_u$ which is the inverse image of $V$
under $u\mto t\cdot u$, one therefore sees that for every element $x$
of $(T\cdot W)(k)$, there exists a conjugate of $H$ which contains $T$
and $x$. As one has already noted, such a conjugate is of the form
$\operatorname{int}(g)\cdot H$, where $g\in N(T)(k)$. Consider then the
morphism
\[
 f:N(T)\times H\to G
\]
defined by $f(g,h)=\operatorname{int}(g)\cdot h=ghg^{-1}$; then the image
of $f$ contains $T\cdot W$, so, since $N(T)$ is a finite union of
translates $C\cdot g_i$ (where $g_i\in N(T)(k)$), as $C$ has finite
index in $N(T)$, it follows that there exists a dense open subset $V'$
of $C=T\cdot C_u$ which is contained in the image of $(C\cdot g_i)\times H$
under $f$. Replacing $H$ by $\operatorname{int}(g_i)\cdot H$, one may
suppose $g_i=e$, i.e. $f(C\times H)\supset V'$. Thus for every
$u\in V'(k)$, there exist $v\in C(k)$ and $h\in H(k)$ such that
\[
 v^{-1}hv=u\qquad\text{whence }vuv^{-1}\in H(k),
\]
whence, setting
\[
 C'=C\cap H=\Centr_H(T),
\]
$vuv^{-1}\in C'$, whence $\operatorname{int}(v)\cdot C'\to u$.
This proves that the union of the conjugates of $C'$ in $C$ (by elements
of $C(k)$) is dense, which implies (as one has seen for the pair $(G,H)$
in place of $(C,C')$) that $C'$ has finite index in its normalizer in $C$.
By BIBLE 7, the lemma of \No 1, it follows that $C'=C$, hence $C=H$, which
proves (vi) $\Rightarrow$ (i) when $G$ is affine.
% typo?: kept as printed: int(v) C' -> u, and C=H.

In the general case, we proceed by induction on $n=\dim G$, the assertion
being trivial if $n=0$. Let $Z$ be the center of $G$, and distinguish
two cases:

\noindent $1^\circ)$ $\dim Z\cap H>0$; then, setting $G'=G/Z\cap H$, one
has $\dim G'<n$; on the other hand, hypothesis (vi) on $H$ implies the
same condition for the image $H'$ of $H$ in $G'$, so $H'$ contains a
Cartan subgroup $C'$ of $G'$, hence $H$ contains the inverse image $C$
of $C'$, which is a Cartan subgroup by XII 6.6 e).

\noindent $2^\circ)$ $\dim Z\cap H=0$, so the canonical morphism $H\to G/Z$
is a finite morphism, and since $G/Z$ is affine by XII 6.1, it follows
that $H$ is affine, hence every homomorphism from $H$ to an abelian
variety is zero (and even every morphism of preschemes from $H$ to an
abelian variety is zero): this follows from the fact that a smooth
connected affine algebraic group over an algebraically closed field
is a rational variety, or merely that it is the union of its Borel
subgroups (BIBLE 6 th. 5 b)); now it follows very easily from the
structure theorems, BIBLE 6.2 and 6.3, that a smooth connected solvable
affine group is a rational variety. Let us now use Chevalley's structure
theorem for $G$, according to which $G$ is an extension of an abelian
variety $A$ by a smooth affine group. Then the image of $H$ in $A$ is
zero, $H$ being affine; on the other hand it is identical to $A$, since
the union of its conjugates in $A$ must be dense. Thus $A=0$, so $G$ is
affine, and one is reduced to the case already treated. This completes
the proof of 2.1.
\end{proof}

\begin{corollary}{2.3}\label{XIII.2.3}
Suppose that the equivalent conditions of 2.1 hold.

\noindent a) Let $k(X)$ (resp. $k(G)$) be the field of rational functions
of $X$ (resp. $G$); then $k(X)$ is a finite separable extension of $k(G)$.
Denote its degree by $d$.

\noindent b) Let $T$ be a maximal torus of $G$ contained in $H$ (which
exists by form 2.1 (ii)), and let $C$ be the corresponding Cartan subgroup
of $G$. Then $C\subset H$. On the other hand, $\Norm_G(T)$ is a smooth
subgroup of $G$ and
$\Norm_G(T)\cap\Norm_G(H)=\Norm_{\Norm_G(H)}(T)$ is a smooth subgroup of
it of finite index equal to $d$ (defined in a)). The number of conjugates
of $H$ containing a given maximal torus or a given Cartan subgroup is
equal to $d$.

\noindent c) Let $U$ be the largest open subset of $G$ such that
$\psi:X\to G$ induces a morphism $\psi^{-1}(U)\to U$ which is finite and
\'etale. Then $U$ is a dense open subset, and for $g\in G(k)$, one has
$g\in U(k)$ if and only if there exist exactly $d$ conjugates of $H$
containing $g$, or again if and only if there exist at least $d$ distinct
conjugates $H_i$ of $H$ containing $g$ such that for every $i$,
$(\mathfrak g/\mathfrak h_i)^{\ad(g)}=0$ (where $\mathfrak h_i=\Lie(H_i)$).
\end{corollary}
Assertion a) follows from the fact that $\psi$ is generically \'etale
(which was stated at the end of 2.1); this also implies that the open
subset $U$ introduced in c) is nonempty, i.e. dense, and the two stated
characterizations of the elements of $U(k)$ (taking into account that
$\psi$ is separated, $X$ integral and $G$ integral and normal, SGA1 I
10.11, and the fact that $(\mathfrak g/\mathfrak h_i)^{\ad(g)}=0$ means
that $\psi$ is \'etale at the point $x_i$ of $\psi^{-1}(g)$ corresponding
to $H_i$). If $H$ contains the maximal torus $T$ of $G$, then the
centralizers of $T$ in $H$ and $G$ have the same dimension and are smooth
and connected (XII 6.6 b)), hence are equal, which proves that $C\subset G$.
% typo?: kept as printed: C subset G in the proof of b).
Moreover, one knows that the normalizer of $T$ in a smooth group which
contains it is smooth (XI, 2.4 bis), so $\Norm_G(T)$ and
$\Norm_{\Norm_G(H)}(T)$ are smooth (N.B. it was noted that $N=\Norm_G(H)$
is smooth, at the end of statement 2.1); moreover, $\Norm_N(T)$ contains
$C$, which has finite index in $N(T)$, hence it has finite index in
$N(T)$. Using the conjugacy theorem for maximal tori of $H$, one sees
that the index in question is equal to the number of conjugates of $H$
which contain $T$, or, what amounts to the same thing, which contain $C$.
Now, since the union of the conjugates of $C$ in $G$ is dense (by
2.1 (i) $\Rightarrow$ (vi) applied to $C$ in place of $H$), and the open
subset $U$ defined in c) is clearly stable under inner automorphisms,
one sees that $C\cap U\ne\varnothing$. Proceeding as in the proof of
implication (vi) $\Rightarrow$ (ii) of 2.1, one concludes that (making,
if necessary, a harmless change of base field) there exists a
$g\in(C\cap U)(k)$ such that every conjugate of $H$ which contains $g$
contains $T$, and consequently also $C$. Thus the conjugates of $H$
containing $C$ are those containing $g$, and since $g\in U(k)$, their
number is equal to $d$, which finishes establishing b). One will note
that we have in fact established that the set of conjugates of $H$
containing $T$ is a homogeneous set under the group of rational points of
\[
 W_G(T)=\Norm_G(T)/\Centr_G(T),
\]
which proves in particular that
\[
 d\leq\text{order of the Weyl group of }G.
\]

\begin{corollary}{2.4}\label{XIII.2.4}
With the notation of 2.1, the following conditions are equivalent:
\begin{enumeratei}
\item $\psi:X\to G$ is a birational morphism.
\item There is one and only one conjugate of $H$ containing a given
Cartan subgroup of $G$.
\item $H$ contains a Cartan subgroup $C$ of $G$, and
$\Norm_G(H)\supset\Norm_G(C)$.
\item There exists a nonempty open subset $V$ of $G$ such that
$g\in V(k)$ implies that $g$ is contained in exactly one conjugate of $H$.
\end{enumeratei}
\end{corollary}
This is clear by 2.1 and 2.3.

\begin{corollary}{2.5}\label{XIII.2.5}
Suppose that the conditions of 2.4 hold, and let $g\in G(k)$.
The following conditions are equivalent:
\begin{enumeratei}
\item $g\in U(k)$, where $U$ is defined in 2.3 c), i.e. $g$ is contained
in one and only one conjugate of $H$.
\item The set of conjugates of $H$ containing $g$ is finite and nonempty.
\item The scheme $\psi^{-1}(g)$ ``of conjugates of $H$ containing $g$''
contains an isolated point.
\item There exists a conjugate $H'$ of $H$ containing $g$, and one has
$(\mathfrak g/\mathfrak h')^{\ad(g)}=0$, where $\mathfrak h'=\Lie(H')$.
\end{enumeratei}
Finally, $U$ is also the largest open subset of $G$ such that $\psi$
induces an isomorphism $\psi^{-1}(U)\isomto U$.
\end{corollary}
The equivalence of (i), (ii), (iii), as well as the last assertion,
follows from the ``Main Theorem''\nde{Of Zariski!} applied to the
birational morphism $\psi:X\to G$, taking into account that $G$ is normal.
The equivalence of these conditions with (iv) follows at once from the
last assertion of 2.1 characterizing the set of elements of $X$ at which
$\psi$ is \'etale.

\begin{theorem}{2.6}\label{XIII.2.6}
Let $G$ be a smooth connected algebraic group over an algebraically closed
field $k$, $C$ a Cartan subgroup, associated to a maximal torus $T$,
$N=\Norm_G(C)=\Norm_G(T)$ (cf. XII 8.4), and let $X=G\times^N C$, where
$N$ acts on the left factor $G$ by right translations, and on the right
factor $C$ by inner automorphisms; let $\psi:X\to G$ be the canonical
morphism.

\noindent a) The morphism $\psi$ is birational.

\noindent b) Let $U$ be the largest open subset of $G$ such that $\psi$
induces an isomorphism $\psi^{-1}(U)\isomto U$ (cf. 2.5). Let
\[
 \rho=\rho_\nu(G)=\dim C
\]
be the nilpotent rank of $G$. Then for every $g\in G(k)$, the multiplicity
of the eigenvalue $1$ in $\ad(g)$ acting on $\mathfrak g$ is at least
$\rho$, and for it to equal $\rho$, it is necessary and sufficient that
one have $g\in U(k)$.
\end{theorem}
\begin{proof}
Since condition 2.1 (i) is verified, one may apply 2.4 (iii)
$\Rightarrow$ (i), which establishes (a). In BIBLE 7 (in the case where
$G$ is affine) the points of $U(k)$ are called the \emph{regular points}
of $G(k)$, and we shall follow this terminology by calling $U$ the open
subset of regular points of $G$. (N.B. The proof given in BIBLE of the
fact that the set in question is itself open is incorrect, but we have
obtained it in the present expos\'e under more general conditions.)
% typo: ou peut -> on peut; expoee -> expose.

Let us prove b), and for this purpose introduce for every $g\in G(k)$
the characteristic polynomial
\[
 P(\ad(g),t)=t^n+c_1(g)t^{n-1}+\cdots+c_n(g);
\]
one sees at once (replacing $k$ by an arbitrary algebra over $k$) that
the $c_i(g)$ come from well-determined sections
\[
 c_i\in\Gamma(G,\OO_G).
\]
When $g\in G(k)$ is an element contained in a Cartan subgroup (for
example, a regular element), which we may suppose to be $C$, then by
2.5 (iv) one sees that one has $(\mathfrak g/\mathfrak c)^{\ad(g)}=0$
if and only if $g$ is regular (where $\mathfrak c$ denotes the Lie
algebra of $C$); on the other hand, since $C$ is nilpotent one sees
immediately that $\ad_{\mathfrak c}(g)$ has only the eigenvalue $1$,
which proves that the multiplicity of the eigenvalue $1$ in
$\ad_{\mathfrak g}(g)$ is $\geq\rho$, and equal to exactly $\dim C=\rho$
if and only if $g$ is regular. In particular, the above polynomial
is divisible by $(t-1)^\rho$. Since the divisibility relation by
$(t-1)^\rho$ is expressed by linear relations (with integer coefficients)
between the coefficients of the polynomial, and these relations hold
for $g\in U(k)$, $U$ being a dense open subset, it follows ($G$ being
reduced) that they hold for every $g$; in fact one has a relation
\[
 t^n+c_1t^{n-1}+c_0=(t-1)^\rho
 (t^{n-\rho}+b_1t^{n-\rho-1}+\cdots+b_{n-\rho})\tag{$\dagger$}
\]
% typo?: kept as printed: c_0 and missing intermediate terms on the left.
in the polynomial ring over $\Gamma(G,\OO_G)$; in particular for every
$g\in G(K)$, $\ad(g)$ has the eigenvalue $1$ with multiplicity at least
$\rho$. Moreover, one has seen that one has equality if $g$ is regular;
let us prove the converse. For this, suppose first that $G$ is affine,
and write $g$ as the product
\[
 g=g_sg_u
\]
of its semisimple part and its unipotent part (BIBLE 4 \No 4); then
\[
 \ad(g)=\ad(g_s)\ad(g_u)
\]
is the analogous decomposition of $\ad(g)$ (\loccit cor. to th. 3), and
consequently $\ad(g)$ and $\ad(g_s)$ have the same eigenvalues (counted
with their multiplicities); in particular the eigenvalue $1$ occurs
with the same multiplicity in $\ad(g)$ and in $\ad(g_s)$.

On the other hand, by BIBLE 7 th. 2 cor. 1, $g$ is regular if and only
if $g_s$ is. Thus, to prove b), one may suppose $g=g_s$, i.e. $g$
semisimple, hence contained in a maximal torus by BIBLE 6 th. 5 c),
and a fortiori in a Cartan subgroup, a case which has already been
treated. This proves b) in the case where $G$ is affine. In the general
case, let $Z=\Centr(G)_{\mathrm{red}}$; then by XII 6.6 e) the Cartan
subgroups of $G$ are the inverse images of those of $G'=G/Z$, so $g$
is regular in $G$ if and only if its image $g'$ in $G'$ is regular in
$G'$. On the other hand, since $Z$ is smooth, the Lie algebra
$\mathfrak g'$ of $G'$ is precisely $\mathfrak g/\mathfrak z$, where
$\mathfrak z=\Lie(Z)$, and $\ad(g')$ is precisely
$\ad(g)_{\mathfrak g/\mathfrak z}$, so the multiplicity of the eigenvalue
$1$ in $\ad(g)$ is equal to $d=\dim Z$ plus the multiplicity of the
eigenvalue $1$ in $\ad(g')$, whence it follows at once that the first
is equal to the nilpotent rank of $G$ if and only if the second is
equal to the nilpotent rank of $G'$. Thus one is reduced to the case
of $G'$; now, $G'$ being affine by XII 6.1, this case has already been
treated. This completes the proof of 2.6.
% typo: on vu -> on a vu.
\end{proof}

\begin{corollary}{2.7}\label{XIII.2.7}
With the notation of the preceding proof\nde{Cf. ($\dagger$) above.}, let
\[
 b=1+b_1+\cdots+b_{n-\rho}\in\Gamma(G,\OO_G).
\]
Then the open subset of regular points of $G$ is given by
\[
 U=G_b
\]
(the set of points of $G$ at which $b$ is invertible); in particular $U$
is an affine open subset if $G$ is affine.
\end{corollary}

% original p. 267
\begin{corollary}{2.8}\label{XIII.2.8}
Let $H$ be a smooth connected algebraic subgroup of $G$ containing a
Cartan subgroup of $G$.

\noindent a) Let $C$ be an algebraic subgroup of $H$. For $C$ to be a
Cartan subgroup of $H$, it is necessary and sufficient that it be a
Cartan subgroup of $G$.

\noindent b) Let $g\in G(k)$, and let $d$ be the integer introduced
in 2.3. For $g$ to be a regular point of $G$, it is necessary and
sufficient that there exist exactly $d$ conjugates $H_i$ of $H$ containing
$g$, and that for each $i$, $g$ be a regular element of $H_i$, or again
that there be at most $d$ conjugates of $H$ containing $g$, and that
$g$ be regular in one of them. If this is so, and if $C$ is the unique
Cartan subgroup of $G$ containing $g$, then the conjugates of $H$
containing $g$ are the conjugates of $H$ containing $C$.

\noindent c) Let $g\in H(k)$; for $g$ to be regular in $G$, it is
necessary and sufficient that it be regular in $H$, and that one have
\[
 (\mathfrak g/\mathfrak h)^{\ad(g)}=0.
\]
\end{corollary}
\begin{proof}
Let us prove a). Under either hypothesis on $C$, the unique maximal
torus $T$ of $C$ is a maximal torus of $H$ and of $G$ ($H$ having the
same reductive rank as $G$), so, since
$\Centr_H(T)\subset\Centr_G(T)$ are smooth connected groups of the
same dimension, they are equal; hence it amounts to the same thing to
say that $C$ is equal to one or the other of these two groups, which
proves a).

Let us prove b). Suppose first that $g$ is regular for $G$, and let
$C$ be the unique Cartan subgroup of $G$ containing $g$; then by 2.3 b)
there exist exactly $d$ conjugates $H_i$ of $H$ containing $C$. Since
$(\mathfrak g/\mathfrak c)^{\ad(g)}=0$, i.e. $\ad(g)$ does not have
the eigenvalue $+1$ on $\mathfrak g/\mathfrak c$, one has a fortiori
$(\mathfrak g/\mathfrak h_i)^{\ad(g)}=0$, so by 2.3 c) there are exactly
$d$ conjugates of $H$ containing $g$, namely the $H_i$. For such an
$H_i$, a Cartan subgroup of $H_i$ containing $g$ is a Cartan subgroup
of $G$ containing $g$, by a), and is therefore equal to $C$, which
proves that $g$ is regular in $H_i$. Conversely, suppose that there
exist at most $d$ conjugates $H_i$ of $H$ containing $g$, and that $g$
is regular in one of them, which may be supposed equal to $H$.
Let us prove that $g$ is regular in $G$. Since $g$ is regular in $H$,
it is contained in a unique Cartan subgroup $C$ of $H$; by a), this
is a Cartan subgroup of $G$. Let $C'$ be a Cartan subgroup of $G$
containing $g$; let us prove $C'=C$ (which will prove that $g$ is
regular in $C$).
% typo?: kept as printed: regular in C rather than G.
Indeed, by 2.3 b), there exist exactly $d$ conjugates of $H$ containing
$C'$, and since these contain $g$, they are necessarily the $H_i$,
so the $H_i$, and in particular $H$, contain $C'$.
% The proof of Corollary 2.8 continues in en-03.tex.
